\chapter{Zusammenfassung und zukünftige Arbeiten}\label{ch:Zusammenfassung und zukünftige Arbeiten}

\lipsum[15-16] % remove

\controlquestions{
Dieses Kapitel soll die wichtigsten Erkenntnisse klar und prägnant zusammenfassen und einen Ausblick auf mögliche Weiterentwicklungen oder offene Fragestellungen geben.
\begin{itemize}
    \item Fazit:
        \begin{enumerate}
            \item Wird die zentrale Forschungsfrage noch einmal klar benannt?
            \item Sind die wichtigsten Ergebnisse und Erkenntnisse verständlich und prägnant zusammengefasst – ohne Wiederholung ganzer Passagen?
            \item Wird nachvollziehbar gezeigt, inwiefern die Ziele der Arbeit erreicht wurden?
            \item Wird klar, wie die Arbeit zum Stand der Forschung und/oder Praxis beiträgt?
            \item Werden keine neuen Inhalte oder Details eingeführt?
            \item Ist die Zusammenfassung kompakt, aber dennoch inhaltlich vollständig?
            \item Wird ein klarer, reflektierter letzter Eindruck hinterlassen?
            \item Sind alle Aspekte des Abstracts und des Kapitels “Motivation” gespiegelt, d.\,h.~alle dort erwähnten Themen/Punkte werden nun aus dem Blickpunkt der abgeschlossenen Arbeit erneut dargestellt?
        \end{enumerate}
    \item Zukünftige Arbeiten:
        \begin{enumerate}
            \item Wurden im Ausblick die (wichtigsten) sich neu ergebenen wissenschaftlichen Fragestellungen besprochen?
            \item Wurden realistische und relevante Vorschläge für weiterführende Arbeiten gemacht?
            \item Wurde begründet, warum bestimmte Aspekte oder Fragestellungen in künftiger Forschung weiterverfolgt werden sollten?
            \item Wurde reflektiert, welche Limitationen der Arbeit zu weiterem Forschungsbedarf führen?
            \item Sind die vorgeschlagenen zukünftigen Arbeiten konkret (z.\,B.~methodisch, thematisch oder technisch eingegrenzt)?
        \end{enumerate}
\end{itemize}

}
